\chapter{第二十一讲\quad 破坏经济秩序犯罪}
\label{lecture:21}
\noindent{\small\hyperref[ref:lecture:21]{查阅本讲扩展规则}}\par

本讲按教材八节展开。以 \(B_i\) 表示第 \(i\) 条所选分支的构成要件满足，\(Q_i\) 表示该分支规定的数额、危险、损害或情节条件满足。两个罪名中的 \(Q\) 即使写法相同，也必须分别核对。以下各例假定明示的事实已经查清，未展开的刑事责任一般条件另行审查。

\section{第一节\quad 生产、销售伪劣商品罪}

\paragraph{法条与定义}
令 \(S\) 表示销售额，\(F\) 表示掺杂掺假、以假充真、以次充好或以不合格产品冒充合格产品，\(P\) 表示生产或销售。第140条既遂的典型分支为
\[
P\land F\land (S\ge 5\text{万元})\ \Longrightarrow\ B_{140}.
\]
这条式子描述销售额达标的分支。未销售的案件要另查货值及未遂规则。食品药品各罪的入罪谓词另设。第143条的食品须达到足以造成严重食物中毒事故或者其他严重食源性疾病的危险 \(D\)，第144条有毒、有害食品分支以掺入有毒、有害非食品原料或销售明知掺有该原料的食品为核心，不以第143条的 \(D\) 为附加条件。第141条假药、第142条劣药与第142条之一妨害药品管理各用法定对象和行为，后两条分别另查严重健康危害和法定健康危险等条件。

\paragraph{推导}
甲为出售在食品中加入有毒非食品原料，已经卖出六万元。案情给出 \(P,F,S\ge5\text{万元}\)，推出 \(B_{140}\)。原料性质与生产行为又推出 \(B_{144}\)。第149条第二款规定同时构成时依处罚较重的规定定罪处罚，因此对两个候选分别确定适用法定刑，再择重一罪。
\[
B_{140}\land B_{144}
\xrightarrow{\text{第149条第二款}}
\mathrm{SelectHeavier}(140,144).
\]
\(\mathrm{SelectHeavier}\) 表示适用较重规定，不把两个法定刑相加。

\paragraph{边界}
换成一般质量不合格的食品，且证实达不到第143条规定的危险，便不能由「是食品」推出该条。若仍有 \(F\) 和达标销售额，第149条第一款保留第140条的入口。劣药与假药的对象和入罪条件也不能互换。

还要区分药品资格和实际危险。甲销售的药品无相应批准证明，但该药在境外合法上市且经证据查明不符合假药的认定标准，也没有第142条之一所要求的严重健康危险。则「无批准」不能同时充当假药谓词和危险谓词。须分别核对药品管理规范与其他可能适用的条款。

\section{第二节\quad 走私犯罪}

\paragraph{法条与模型}
记 \(C\) 为故意逃避海关监管，\(O\) 为普通货物、物品，\(D_s\) 为相应条款规定的特殊对象及进出境方向相合，\(T\) 为偷逃应缴税额。教材把案件先按对象与方向分流。第153条的普通货物分支用
\[
C\land O\land\bigl(Q_T(T)\lor R\bigr)\Longrightarrow B_{153},
\]
其中 \(Q_T\) 表示税额达到适用档，\(R\) 表示一年内曾因走私受过二次行政处罚后又走私。特殊对象则用 \(C\land D_s\land Q_s\) 核对第151、152条等规定，不能借用普通货物的税额入口。第155条列举按走私罪论处的特定行为，第156条处理与走私者通谋而提供便利的情形。

\paragraph{推导}
甲藏匿应申报的普通进口手表，绕过海关，税额经核定达到入罪档。由隐匿绕监管及故意得 \(C\)，由手表性质得 \(O\)，由核定得 \(Q_T\)，依第153条推出 \(B_{153}\)。乙事先商定为该次走私提供运输，若已证明通谋和帮助，则另按第156条审查乙的共犯责任。

\paragraph{边界}
给甲改成运出国家禁止出口的文物，且对象和方向均经核实，\(O\) 的分支失效，应查第151条。再换成依法允许出口的普通文物，不能保留原来的 \(D_s\)。事后才知情的普通运输也不能倒推出事先通谋。

\paragraph{对象错误：教材的包容评价}
令 \(K_s\) 为对所指具体特殊对象的认识，\(K_g\) 为对所走私物品属于禁限对象的概括认识，二者强度不同。教材处理具体对象认错时，在已证明客观对象符合一般禁止进出口分支的前提下，进一步审查是否具有 \(K_g\)。例如甲以为箱内是假币，实为禁止进口的淫秽物品。即使概括认识有证据，仍不能据此证明甲具有第152条第一款所要求的牟利或传播目的。能否落到一般走私分支，要另核该条的客观兜底范围及故意。这个模型讨论构成要件重合，不能当作三个走私条文的适用范围互为集合包含的证明。

\section{第三节\quad 妨害对公司、企业的管理秩序罪}

\paragraph{四项参照}
法释〔2026〕6号第8条将第163、164、271、272条的定罪量刑标准分别联系到受贿、行贿（单位行贿）、贪污、挪用公款的标准。这是对入罪和量刑标准的参照，原有各罪的主体和行为仍须分别核对。决定是否追责和量刑时，还应综合犯罪性质、情节和社会危害性。

\paragraph{法条与模型}
分别记录身份 \(I\)、本职便利 \(V\)、收财或给财 \(G\)、为他人谋利 \(M\) 或谋取不正当利益 \(U\)。第163条的基本受贿分支核对非国家工作人员身份、利用职务便利、索取或者非法收受财物、为他人谋利及入罪数额。第164条给财分支核对为谋取不正当利益而给公司、企业或其他单位工作人员财物以及数额条件。第165、166条经刑法修正案（十二）修改。以第165条为例，国有公司、企业董事、监事、高管的分支核对经营同类营业并取得达到规定数额的非法利益。其他公司、企业相应人员的分支须另核对违反法律、行政法规及公司、企业利益重大损失。

\paragraph{推导}
私营公司采购经理利用选商权限索取财物，并承诺给供货商订单。身份为 \(I_{\mathrm{private}}\)，选商权限推出 \(V\)，索财及承诺推出该案的收财行为与 \(M\)。若数额已达标，依第163条相应分支得 \(B_{163}\)。供货商是否成立第164条，还须独立证明给财、\(U\) 与数额。即使买卖双方都在场，也不能由一方成立直接推出另一方成立。

\paragraph{边界}
私营公司董事未经许可兼职，不足以证明第165条。即使已经证明利用职务便利经营同类营业，仍须核对其所属身份分支的法规违反及重大损失。第166条为亲友非法牟利还应检查法定交易或业务安排类型和相应损失。

\section{第四节\quad 破坏金融管理秩序罪}

\paragraph{法条与模型}
本节分三组输入。货币组以真币是否被加工为起点区分第170条伪造与第173条变造，再区分第171条出售、购买、运输和第172条持有、使用。伪造者又出售、运输自制假币，依第171条第三款按伪造货币罪从重。洗钱组设 \(U_7\) 为第191条列举的七类上游犯罪事实成立，\(W\) 为洗钱行为，\(H\) 为掩饰、隐瞒来源和性质的目的。贷款组区分高利转贷的套取信贷资金、高利转贷牟利、违法所得条件（第175条），骗取贷款等致金融机构重大损失（第175条之一），和金融机构工作人员违法发放贷款、数额巨大或者造成重大损失（第186条）。第177条核对金融票证伪造、变造，第177条之一另核对信用卡管理及信用卡信息。

\paragraph{推导}
甲受贿得款，另虚构交易，把款项转成店铺营业收入，目的是掩饰款项来源。已证实受贿事实，故有 \(U_7\)。虚构交易转换资金与该目的给出 \(W,H\)。依第191条及洗钱解释（法释〔2024〕10号）第1、5条，推出自洗钱分支的 \(B_{191}\)。上游受贿和后续洗钱分别评价。上游案件虽尚未生效裁判，有证据证明确实存在仍可认定上游事实（同解释第7条）。

\paragraph{边界}
改变上游为单纯盗窃，则失去 \(U_7\)，第191条不能沿该式推出，另按第312条及相关解释评价适格行为。普通生活消费也不能仅因钱来自受贿即推出 \(H\land W\)。

\section{第五节\quad 金融诈骗罪}

\paragraph{法条与定义}
令 \(J\) 表示非法占有目的，\(K\) 表示真实信用卡，\(A\) 表示冒用，\(Z\) 表示盗窃，\(U\) 表示使用。第196条的冒用分支与第三款的特别规则分两步检查。
\[
\begin{aligned}
A\land K\land U\land Q_{196}&\Longrightarrow B_{196}^{\text{候选}},\\
Z\land K\land U&\xrightarrow{\text{第196条第三款}}\text{按盗窃罪规则评价}.
\end{aligned}
\]
冒用可以涉及拾得他人卡后使用、骗取他人卡后使用，或者非法取得他人信用卡信息资料并通过互联网、通讯终端等使用。恶意透支另用持卡人身份、非法占有目的 \(J\)、超过规定限额或期限 \(E\)、两次有效催收 \(C_2\) 和再次有效催收后超过三个月仍不归还 \(R_3\)。根据2018年修正的信用卡刑事司法解释第6至9条，\(C_2\) 包含两次催收至少间隔三十日等条件。入罪数额看刑事立案时未还实际透支本金，不加计利息和费用。即使 \(E,C_2,R_3\) 均成立，仍须另证 \(J\)。

\paragraph{推导}
甲秘密窃取乙的真实银行卡，取得密码后在自动取款机取出款项，金额已达盗窃入罪条件。第一步由取得方式得到 \(Z\)，由发卡资料确认 \(K\)。第二步由取款日志及实际取现得到 \(U\) 和既遂取财事实。第三步依第196条第三款落在盗窃特别规则，得盗窃罪的结论。若只改成拾得卡后冒名取款，且证明冒用及达标，\(Z\) 消失，改走信用卡诈骗的冒用分支。

\paragraph{相邻模型与边界}
第193条另核对以非法占有为目的诈骗银行或其他金融机构贷款及数额，第198条核对特定主体、保险骗取类型与数额，第192条核对 \(J\land\) 诈骗方法非法集资及数额。第176条考查非法吸收公众存款或变相吸收公众存款、扰乱金融秩序，不能把还款不能直接写成 \(J\)。同理，非法取得支付软件账号与直接使用银行卡信息不是同一个事实，先辨实际动用的是支付余额、银行结算资金还是信贷资金，再选规则。

\section{第六节\quad 危害税收征管罪}

\paragraph{法条及司法解释}
设某一已确定口径内应纳税额 \(T>0\)，逃税额 \(E\)，比例 \(\rho=E/T\)。令 \(D\) 为采取欺骗、隐瞒手段虚假申报或者不申报。第201条第一款的基础分支可写为
\[
D\land (E\ge 10\text{万元})\land(\rho\ge 10\%)
 \Longrightarrow B_{201}^{\mathrm{base}} .
\]
十万元来自法释〔2024〕4号第2条，比例来自刑法第201条。适用比例前应核对同解释第4条的期间、税种及口径。第一款的行政处理后不追究规则需要税务机关依法下达追缴通知、补缴税款、缴纳滞纳金、已受行政处罚及解释所定程序条件。五年内因逃避缴纳税款受过刑事处罚或者被税务机关给予二次以上行政处罚的，不能走这一例外。扣缴义务人逃避已扣已收税款用第二款，不能套用第一款的比例与第四款的特别规则。

\paragraph{推导}
纳税人以两套账隐匿收入。以同一已核定口径，\(T=100\) 万元，\(E=12\) 万元，则 \(12/100=12\%\)。手段与两个阈值都成立，推出 \(B_{201}^{\mathrm{base}}\)。若另查清已依法按期完成补缴、缴滞纳金、履行行政处罚且无例外限制，则依第四款得不追究刑事责任。把 \(T\) 改为 \(150\) 万元，\(E\) 不变，\(\rho=8\%\)，基础分支即不能推出。

\paragraph{相邻模型与边界}
第204条骗取出口退税和第201条的分流关注是否超出已缴税款范围。对已缴 \(T_0\) 而以欺骗方法骗回 \(R>0\) 的案件，第204条第二款将金额拆成 \(\min(R,T_0)\) 和 \(\max(0,R-T_0)\)，前段按第201条，超出部分按第204条第一款分别审查。第205条核对抵扣、骗取退税功能与法定虚开行为，第205条之一核对其他发票虚开及严重情节。按税收司法解释第10条，因虚增业绩、融资、贷款等不以骗抵税款为目的，且没有因抵扣造成税款被骗损失的虚开，不按第205条定罪，构成其他犯罪另行评价。发票记录不实本身不能代替完整审查。

\section{第七节\quad 侵犯知识产权罪}

\paragraph{法条与模型}
第217条设 \(P\) 为营利目的，\(N\) 为未经权利人许可，\(V\) 为法定侵权方式，\(Q\) 为违法所得较大或其他严重情节。典型侵权分支为 \(P\land N\land V\land Q\Rightarrow B_{217}\)。第213条另核对同一种商品、服务与相同注册商标及情节。第219条先证实商业秘密的非公知性、商业价值及合理保密措施，再核对法定获取、披露、使用等行为与严重情节。

\paragraph{推导}
甲搭建收费网站，把未获许可的电影放在可由公众选定时间和地点观看的页面，因该行为违法所得四万元。收费获利证明 \(P\)，权利人和授权链证实 \(N\)，依2025年知识产权刑事司法解释第12条得到 \(V\)。四万元达到同解释第13条的三万元违法所得门槛，得 \(Q\)。合取后依第217条得到相应分支。

\paragraph{边界}
权利人已有效许可该种传播，\(N\) 不成立，即使四万元也不能推出该分支。对于商标，视觉近似在民事上可能有关，不能直接代替刑法及其司法解释中的「相同」认定。公开流传且缺乏保密条件的资料不能仅因标注「内部」便当作商业秘密。

\section{第八节\quad 扰乱市场秩序罪}

\paragraph{法条与模型}
第225条的三个入口是违反国家规定 \(N\)、落入法定非法经营类型 \(A\)、严重扰乱市场秩序 \(Q\)。第226条关注暴力、威胁与法定强迫交易行为及情节。第229条依中介职责、故意提供虚假证明文件或严重不负责任致重大失实等分支分别查证。第224条关注签订、履行合同中的非法占有目的、骗取对方财物和数额。第224条之一关注以经营等活动为名、缴费或购物等取得资格、层级及以发展人数计酬返利，组织领导者引诱、胁迫发展、骗取财物并扰乱经济社会秩序。

\paragraph{推导}
甲未获许可，经营国家规定专卖的商品。证据确认其长期批量经营并已满足情节严重标准。许可义务和无许可给出 \(N\)，商品性质和经营给出第225条第1项的 \(A\)，量与经营情况给出 \(Q\)，合并推出该项分支。若甲仅有经营手续瑕疵而未落入任何法定类型，\(N\) 即使存在，也不能推出 \(B_{225}\)。

\paragraph{边界}
合同一方因市场变化无力履行，不能由违约推出非法占有目的。若题干已证实签约时虚构履约能力并以占有为目的骗取达标款项，再核对第224条具体方式。考查中介报告时应先定位行业职责，再分别查故意与过失，不能用「报告结论错了」代替这两项。

\paragraph{闭卷记忆}
\begin{enumerate}
\item 先画八格：货、关、职、金、骗、税、权、市。每格写一个最先检查的输入，依次为对象、进出境方向、身份、金融行为、非法占有目的、税额口径、权利许可、国家规定。
\item 每格再写「动作、阈值、例外」。把 \(E/T\) 写在税格，把 \(U_7,W,H\) 写在金融格，把盗窃后使用真卡的特别规则写在诈骗格。
\item 遇到竞合先列两条可成立的构成要件，最后写明选择规则的条号。
\end{enumerate}
口诀「\textbf{货关职金，骗税权市。先分对象，再核条件}」逐格对应上面八个输入。

\paragraph{闭卷自测与答案}
\begin{enumerate}
\item 普通纳税人同一口径逃税 \(12\) 万元，应纳 \(150\) 万元。假设只有该期该项事实，能否由第201条第一款的数额比例分支入罪？答：\(12/150=8\%\)，金额条件成立而比例条件不成立，该分支推不出。
\item 某人盗窃他人真卡并取现，数额达标。只将取得方式改成拾得，其他条件相同，哪一个谓词变了？答：\(Z\) 从真变假，失去第196条第三款特别规则。若已证实冒用和相应数额条件，另循信用卡诈骗分支。
\item 无证经营何时可推出第225条？答：单凭无证不能。还须有国家规定、对应的法定类型以及严重扰乱市场秩序所需的情节等事实。
\end{enumerate}
